% Part: normal-modal-logic
% Chapter: frame-correspondence
% Section: standard-translation

\documentclass[../../../include/open-logic-section]{subfiles}

\begin{document}

\olfileid{nml}{frd}{st}

\olsection{ద్వితీయ-స్థాయి నిర్వచనీయత}

మోడల్ సూత్రాలతో నిర్వచించగల ప్రతి చట్ర ధర్మమూ
మొదటిస్థాయి నిర్వచనీయమైనది కాదు. అయితే ఏకస్థానిక విధేయాలపై
పరిమాణీకరణను (అంటే, ఏకస్థానిక ద్వితీయ-స్థాయి పరిమాణీకరణను)
అనుమతిస్తే, మోడల్ సూత్రాలతో నిర్వచించగల అన్ని చట్ర ధర్మాలనూ
నిర్వచించవచ్చు. ఒక లోకం వద్ద మోడల్
\tetoken{సూత్రం}{!!{formula}} సత్యమయ్యే షరతులకూ
మొదటిస్థాయి \tetoken{సూత్రాలకూ}{!!{formula}s} మధ్య ఉన్న
క్రమబద్ధమైన సంబంధాన్ని ఉపయోగించడమే ఉపాయం. దీనినే మోడల్
సూత్రాల \emph{ప్రామాణిక అనువాదం} అంటారు. ప్రాప్యత సంబంధం
కోసం ద్విస్థాన \tetoken{విధేయ సంకేతం}{!!{predicate}}~$Q$
మాత్రమే కాక, $!A$లో వచ్చే
\tetoken{ప్రతిజ్ఞావాక్య చరాలు}{!!{propositional variable}s}~$\Obj p_i$
కోసం ఏకస్థానిక \tetoken{విధేయ సంకేతం}{!!{predicate}}~$P_i$
కూడా ఉన్న భాషలోని మొదటిస్థాయి సూత్రాలుగా ఈ అనువాదం
మోడల్ సూత్రాలను మారుస్తుంది.

\begin{defn}
  \emph{ప్రామాణిక అనువాదం}~$\ST_x(!A)$ను ఆగమన
  పద్ధతిలో కిందిలా నిర్వచిస్తాం:
  \begin{enumerate}
  \tagitem{prvFalse}{\indcase{!A}{\lfalse}{$\ST_x(\indfrm) = \lfalse$.}}{}
  \tagitem{prvTrue}{\indcase{!A}{\ltrue}{$\ST_x(\indfrm) = \ltrue$.}%
    \sourcecorrection{OLTENMLFRDST-001}{మూలంలోని సత్య
      విభాగపు సందర్భంలో పొరపాటున అసత్య స్థిరాంకం ఉంది;
      కుడివైపు సమీకరణానికి, పూర్వ సత్య నిర్వచనానికి అనుగుణంగా
      దాన్ని సత్య స్థిరాంకంగా సరిచేశాం.}}{}
  \item \indcase{!A}{\Obj p_i}{$\ST_x(\indfrm) = \Atom{P_i}{x}$.}
  \tagitem{prvNot}{\indcase{!A}{\lnot !B}{$\ST_x(\indfrm) = \lnot \ST_x(!B)$.}}{}
  \tagitem{prvAnd}{\indcase{!A}{(!B \land !C)}{$\ST_x(\indfrm) = (\ST_x(!B)
      \land \ST_x(!C))$.}}{}
  \tagitem{prvOr}{\indcase{!A}{(!B \lor !C)}{$\ST_x(\indfrm) = (\ST_x(!B)
      \lor \ST_x(!C))$.}}{}
  \tagitem{prvIf}{\indcase{!A}{(!B \lif !C)}{$\ST_x(\indfrm) = (\ST_x(!B)
      \lif \ST_x(!C))$.}}{}
  \tagitem{prvIff}{\indcase{!A}{(!B \liff !C)}{$\ST_x(\indfrm) = (\ST_x(!B)
      \liff \ST_x(!C))$.}}{}
  \sourcecorrection{OLTENMLFRDST-002}{ద్విసోపాధిక శాఖలో ఆంగ్ల మూలం గణితపు ముగింపు గుర్తు కంటే ముందే ఆగమన శాఖ వాదనను మూసింది. సూత్రాన్ని మార్చకుండా, తెలుగు లక్ష్యంలో గణితపు ముగింపు గుర్తును ఆ వాదనలోకి మార్చాం.}
  \tagitem{prvBox}{\indcase{!A}{\Box !B}{$\ST_x(\indfrm) =
      \lforall[y][(\Atom{Q}{x,y} \lif \ST_y(!B))]$.}}{}
  \tagitem{prvDiamond}{\indcase{!A}{\Diamond !B}{$\ST_x(\indfrm) =
      \lexists[y][(\Atom{Q}{x,y} \land \ST_y(!B))]$.}}{}
  \end{enumerate}
\end{defn}

ఉదాహరణకు, $\ST_x(\Box p \lif p)$ అనేది $\lforall[y][(\Atom{Q}{x,y}
  \lif \Atom{P}{y})] \lif \Atom{P}{x}$. $\ST_x(!A)$ భాషకు
ఏ \tetoken{నిర్మాణానికైనా}{!!{structure}} ఒక వ్యక్తి క్షేత్రం,
$Q$కు నిర్దేశించిన ద్విస్థాన సంబంధం, ఏకస్థానిక
\tetoken{విధేయ సంకేతాలు}{!!{predicate}s}~$P_i$కు
నిర్దేశించిన ఆ క్షేత్రపు ఉపసమితులు కావాలి. అంటే అటువంటి
\tetoken{నిర్మాణం}{!!a{structure}}లోని అంశాలు
$!A$కు ఒక నమూనాలోని అంశాలే: వ్యక్తి క్షేత్రం లోకాల సమితి;
$Q$కు నిర్దేశించిన ద్విస్థాన సంబంధం ప్రాప్యత సంబంధం;
$P_i$కు నిర్దేశించిన ఉపసమితులు $V(\Obj p_i)$ అనే
కేటాయింపులు. అందువల్ల మోడల్ నమూనాలో $!A$ సంతృప్తికీ,
దానికి అనురూపమైన \tetoken{నిర్మాణంలో}{!!{structure}}
$\ST_x(!A)$ సంతృప్తికీ పొంతన ఉండటం సహజమే:

\begin{prop}\ollabel{prop:st}
  $\mModel{M} = \tuple{W, R, V}$ అనుకుందాం. $\Domain{M'} = W$,
  $\Assign{Q}{M'} = R$, $\Assign{P_i}{M'} = V(\Obj p_i)$ గల
  మొదటిస్థాయి \tetoken{నిర్మాణం}{!!{structure}}~$\Struct{M'}$ను
  తీసుకొని, $s(x) = w$ అనుకుందాం. అప్పుడు
  \[
  \mSat{M}{!A}[w] \text{ iff } \Sat{M'}{\ST_x(!A)}[s]
  \]
\end{prop}

\begin{proof}
  $!A$ నిర్మాణంపై ఆగమనంతో నిరూపించవచ్చు.
\end{proof}

\begin{prop}
  $!A$ ఒక మోడల్ \tetoken{సూత్రం}{!!{formula}} అనీ,
  $\mModel{F} = \tuple{W, R}$ ఒక చట్రం అనీ అనుకుందాం.
  $\Domain{F'} = W$, $\Assign{Q}{F'} = R$ గల మొదటిస్థాయి
  \tetoken{నిర్మాణం}{!!{structure}}~$\Struct{F'}$ను తీసుకుందాం.
  $!A'$ అనే ద్వితీయ-స్థాయి సూత్రం కిందిది అనుకుందాం:
  \[
  \lforall[X_1][\dots\lforall[X_n][\lforall[x][
        \SSubst{\ST_x(!A)}{\subst{X_1}{P_1}, \dots,
          \subst{X_n}{P_n}}]]],
  \]
  ఇక్కడ $P_1$, \dots, $P_n$ అన్నీ $\ST_x(!A)$లోని
  ఏకస్థానిక \tetoken{విధేయ సంకేతాలు}{!!{predicate}s}.
  అప్పుడు
  \[
  \mModel{F} \Entails !A \text{ iff } \Sat{F'}{!A'}
  \]
\end{prop}

\begin{proof}
  $\Sat{F'}{!A'}$ అయ్యేది, $i = 1$, \dots,~$n$కు
  $\Assign{P_i}{M'} \subseteq W$ అయ్యే ప్రతి
  \tetoken{నిర్మాణం}{!!{structure}}~$\mModel{M'}$లోనూ,
  $s(x) \in W$ అయ్యే ప్రతి $s$కూ
  $\Sat{M'}{\ST_x(!A)}[s]$ అయినప్పుడే.
  \olref{prop:st} ప్రకారం, ఇది $\mModel{F}$పై ఆధారపడే
  అన్ని నమూనాలు~$\mModel{M}$లోనూ, ప్రతి లోకం~$w \in W$
  వద్దనూ $\mSat{M}{!A}[w]$ అయితేనే; అంటే
  $\mModel{F} \Entails !A$ అయితేనే.
\end{proof}

\begin{defn}
  చట్రాల వర్గం~$\mClass{F}$ను \emph{ద్వితీయ-స్థాయి
    నిర్వచనీయమైనది} అంటాం, ఒకవేళ ఒక్క ద్విస్థాన
  \tetoken{విధేయ సంకేతం}{!!{predicate}}~$P$ ఉండి,
  ద్వితీయ-స్థాయి పరిమాణీకరణలు ఏకస్థానిక సమితి
  చరాలపైనే ఉండే
  ద్వితీయ-స్థాయి భాషలో ఒక
  \tetoken{వాక్యం}{!!a{sentence}}~$!A$ ఉండి,
  $\mModel{F} = \tuple{W, R} \in \mClass{F}$ అయ్యేది
  $\Domain{M} = W$, $\Assign{P}{M} = R$ గల
  \tetoken{నిర్మాణం}{!!{structure}}~$\Struct{M}$లో
  $\Sat{M}{!A}$ అయినప్పుడే.
\end{defn}

\begin{cor}
  చట్రాల వర్గాన్ని \tetoken{ఒక సూత్రం}{!!a{formula}}~$!A$
  నిర్వచిస్తే, దానికి అనురూపమైన ప్రాప్యత సంబంధాల
  వర్గాన్ని ఏకస్థానిక ద్వితీయ-స్థాయి వాక్యంతో
  నిర్వచించవచ్చు.
\end{cor}

\begin{proof}
  మునుపటి నిరూపణలోని ఏకస్థానిక ద్వితీయ-స్థాయి
  వాక్యం~$!A'$కు కావలసిన ధర్మం ఉంది.
\end{proof}

ఉదాహరణగా, $\Box p \lif p$ అనే
\tetoken{సూత్రాన్ని}{!!{formula}} మళ్లీ చూద్దాం. అది
స్వావర్తనాన్ని నిర్వచిస్తుంది. స్వావర్తనాన్ని
$\lforall[x][\Atom{Q}{x,x}]$ అనే
మొదటిస్థాయి \tetoken{వాక్యం}{!!{sentence}} కూడా
నిర్వచిస్తుంది. అయితే కింది ఏకస్థానిక ద్వితీయ-స్థాయి
\tetoken{వాక్యం}{!!{sentence}} కూడా దాన్ని నిర్వచిస్తుంది:
\[
\lforall[X][\lforall[x][(\lforall[y][(\Atom{Q}{x,y} \lif \Atom{X}{y})]
    \lif \Atom{X}{x})]].
\]
అంటే, ఈ రెండు వాక్యాలూ సమానార్థకాలు. దీన్ని నేరుగా
ఎలా నిర్ధారించవచ్చో చూద్దాం. ముందుగా ద్వితీయ-స్థాయి
వాక్యం \tetoken{ఒక నిర్మాణం}{!!a{structure}}~$\Struct{M}$లో
సత్యమని అనుకుందాం. $x$, $X$ సార్వత్రికంగా
పరిమాణీకరించబడ్డాయి కాబట్టి, మిగిలిన భాగం ప్రతి
$x \in W$కూ ప్రతి సమితి~$X \subseteq W$కూ, ఉదాహరణకు
$R = \Assign{Q}{M}$ అయినప్పుడు
$\Setabs{z}{Rxz}$ అనే సమితికీ, నిలవాలి. కనుక
$s(x) \in W$, $s(X) = \Setabs{z}{Rxz}$ అయ్యే
ప్రతి~$s$కు
$\Sat{M}{\lforall[y][(\Atom{Q}{x,y} \lif
    \Atom{X}{y})] \lif \Atom{X}{x}}$ నిలుస్తుంది.
కానీ $s(X)$ను ఎంచుకున్న విధం వల్ల ఇది
$\Sat{M}{\lforall[y][(\Atom{Q}{x,y} \lif
    \Atom{Q}{x,y})] \lif \Atom{Q}{x,x}}[s]$
అనే అర్థాన్ని ఇస్తుంది. పూర్వభాగం ఎల్లప్పుడూ
చెల్లుబాటు అవుతుంది కాబట్టి అది
$\Atom{Q}{x,x}$కు సమానం. $s(x)$ ఏ లోకమైనా కావచ్చు;
అందువల్ల $\Sat{M}{\lforall[x][\Atom{Q}{x,x}]}$.

ఇప్పుడు $\Sat{M}{\lforall[x][\Atom{Q}{x,x}]}$
అనుకుందాం; $\Sat{M}{\lforall[X][\lforall[x][(\lforall[y][(\Atom{Q}{x,y} \lif
        \Atom{X}{y})] \lif \Atom{X}{x})]]}$ అని చూపుదాం.
ఏదైనా నిర్దేశం~$s$ను తీసుకొని,
$\Sat{M}{\lforall[y][(\Atom{Q}{x,y} \lif
    \Atom{X}{y})]}[s]$ అనుకుందాం. $s'(y)
= s(x)$ అయ్యే $s$ యొక్క $y$-భేదరూపం $s'$ను
తీసుకుంటే $\Sat{M}{\Atom{Q}{x,y} \lif \Atom{X}{y}}[s']$;
అంటే $\Sat{M}{\Atom{Q}{x,x} \lif \Atom{X}{x}}[s]$.
$\Sat{M}{\lforall[x][\Atom{Q}{x,x}]}$ కాబట్టి
పూర్వభాగం సత్యం; అందువల్ల
$\Sat{M}{\Atom{X}{x}}[s]$. చూపవలసింది ఇదే.

నిర్వచించగల కొన్ని చట్రాల వర్గాలు మొదటిస్థాయి
నిర్వచనీయమైనవి కావు. కనుక $!A'$ రూపంలో ఉన్న
ప్రతి ఏకస్థానిక ద్వితీయ-స్థాయి
\tetoken{వాక్యం}{!!{sentence}} మొదటిస్థాయి
\tetoken{వాక్యానికి}{!!{sentence}} సమానార్థకం కాదు.
వాటిలో ఏవి సమానార్థకాలో నిర్ణయించే ప్రభావవంతమైన
పద్ధతి ఏదీ లేదు.

\end{document}
